\documentclass[../../../include/open-logic-section]{subfiles}

\begin{document}

\olfileid{sth}{choice}{banach}
\olsection{బనాక్--టార్స్కీ వైరుధ్యాభాసం}

ప్రతిస్థాపనను సమర్థించడానికి బూలొస్ ప్రయత్నించినట్లే, \emph{బాహ్య} కారణాలను చూపి ఎంపికను కూడా సమర్థించడానికి ప్రయత్నించవచ్చు (\olref[replacement][extrinsic]{sec} చూడండి). ఎంపికను స్వీకరిస్తే ఎన్నో కోరదగిన పరిణామాలు వస్తాయి: ప్రతి కార్డినల్ సంఖ్యను మరొకదానితో పోల్చగలగడం; ప్రతి సమితినీ సుక్రమపరచగలగడం; కార్డినల్ సంఖ్యలను ఒక రకమైన క్రమసంఖ్యలుగా తీసుకోగలగడం; మొదలైనవి.

అయితే ఎంపికకు \emph{అవాంఛనీయ} పరిణామాలున్నాయని కొన్నిసార్లు అంటారు. ముఖ్యంగా \cite{BanachTarski1924} ఫలితమే దీనికి కారణం.

\begin{thm}[బనాక్--టార్స్కీ వైరుధ్యాభాసం ($\ZFC$లో)]
ఏ ఘన గోళాన్నైనా పరిమిత సంఖ్యలో ముక్కలుగా విడదీసి, ఆ ముక్కలను భ్రమణం, స్థానాంతరణం ద్వారా మళ్లీ అమర్చి, అదే ఘన గోళానికి రెండు ప్రతులను ఏర్పరచవచ్చు.
\end{thm}
\noindent
మొదటి చూపులో ఇది ఆశ్చర్యంగా ఉంటుంది. ఆ రెండు ఘన గోళాల మొత్తం ఘనపరిమాణం మొదటి ఘన గోళానికి \emph{రెట్టింపు} అని స్పష్టం. కానీ దృఢ చలనాలు---భ్రమణం, స్థానాంతరణం---ఘనపరిమాణాన్ని మార్చవు. అందువల్ల బనాక్--టార్స్కీ కొత్త పదార్థాన్ని మాయగా సృష్టించినట్లనిపిస్తుంది.

ఇంకా ఆశ్చర్యకరం కూడా ఉంది.\footnote{\citet[Theorem 3.12]{Wagon2016} చూడండి.} ఇలాంటి వాదనతో ఒక బఠాణీ గింజను పరిమిత సంఖ్యలో ముక్కలుగా చేసి, వాటిని భ్రమణం, స్థానాంతరణం ద్వారా మళ్లీ అమర్చి, బిగ్ బెన్ ఆకారం, పరిమాణం గల వస్తువుగా మార్చవచ్చని చూపవచ్చు.

అయితే ఇవేవీ $\ZF$లో మాత్రమే నిరూపితమవవు.\footnote{అయినప్పటికీ, ఎంపికకంటే కచ్చితంగా బలహీనమైన సూత్రాలతో బనాక్--టార్స్కీని నిరూపించవచ్చు; \citet[303]{Wagon2016} చూడండి.} కాబట్టి ఒక నిర్ణయం తీసుకోవాలి: ఎంపికను తిరస్కరించాలా, లేక ఈ ``వైరుధ్యాభాసం''తో జీవించడం నేర్చుకోవాలా?

మేము ఈ ``వైరుధ్యాభాసం''తో జీవించడం నేర్చుకోవాలని సూచిస్తాం. నిజానికి ఇది పెద్ద వైరుధ్యాభాసమే అని మాకు అనిపించదు. ముఖ్యంగా ఈ కింది ఫలితాలకంటే ఇది ఎందుకు ఎక్కువో తక్కువో వైరుధ్యాభాసంగా ఉండాలో మాకు కనిపించదు:\footnote{\citet[276--7]{Potter2004}, \citet[16]{Weston2003}, \citet[31, 308--9]{Wagon2016} ఇతర ఉదాహరణలతో ఇలాంటి విషయాలనే చెబుతారు.}
\begin{enumerate}
	\item $(0,1)$ అంతరంలో ఉన్న బిందువులెన్నో $\Real$లోనూ అన్నే ఉన్నాయి.
	\\\emph{నిరూపణ}: $\tan(\pi(r-\nicefrac{1}{2}))$ను చూడండి.
	\sourcecorrection{OLTESTCHOICEBANACH-001}{మూలంలో ఈ టాంజెంట్ వ్యక్తీకరణకు ఒక అదనపు మూసివేత కుండలీకరణం ఉంది; $(0,1)$ నుంచి వాస్తవ సంఖ్యలకు సరైన ప్రతిచిత్రణగా దాన్ని సరిచేశాం.}
	\item ఒక రేఖలోని బిందువులెన్నో ఒక చతురస్రంలోనూ అన్నే ఉన్నాయి.
	\\\olref[his][set][pathology]{sec}, \olref[his][set][cantorplane]{sec} చూడండి.
	\item స్థలాన్ని నింపే వక్రరేఖలు ఉన్నాయి.
	\\\olref[his][set][pathology]{sec}, \olref[his][set][hilbertcurve]{sec} చూడండి.
\end{enumerate}
ఈ మూడు ఫలితాల్లో దేనికీ ఎంపిక అవసరం లేదు. వాటిని ఇప్పుడు ఆశ్చర్యకరమైన, మనోహరమైన గణిత ఫలితాలుగానే చూస్తాం. బనాక్--టార్స్కీ వైరుధ్యాభాసాన్నీ అలాగే చూడవచ్చేమో.

నిజమే, ఇక్కడ ఒక సాంకేతిక విషయం గమనించాలి; దానికి ప్రశాంతంగా ఆలోచిస్తే చాలు. దృఢ చలనాలు ఘనపరిమాణాన్ని పరిరక్షిస్తాయి. అందువల్ల ఘన గోళాన్ని విడదీసే ఐదు\footnote{వైరుధ్యాభాసాన్ని ``పరిమిత సంఖ్యలో ముక్కలు'' అని చెప్పాం. నిజానికి \citet{Robinson1947} \emph{ఐదు} ముక్కలతో (అంతకంటే తక్కువతో కాదు) ఈ విభజన సాధ్యమని నిరూపించారు. నిరూపణకు \citet[pp.~66--7]{Wagon2016} చూడండి.} ముక్కలన్నీ \emph{కొలవదగినవి} కావు. స్థూలంగా చెప్పాలంటే, విడివిడిగా ఆ ముక్కలకు ఘనపరిమాణం ఇవ్వడానికే అర్థం లేదు. వాటిని చిత్రించలేని, బిందువుల ``అనంతంగా చెదరిన సమూహాలు''గా ఊహించాలి. అలాంటి ``అనంతంగా చెదరిన'' సమితులు విచిత్రంగా అనిపించవచ్చు. కానీ ముందే అందుబాటులో ఉన్న సమితుల \emph{సాధ్యమైన అన్ని} సేకరణలను ఏర్పరచాలనే \stagesacc{} సూత్రం నుంచి వాటి అస్తిత్వం వస్తున్నట్లే కనిపిస్తుంది.

ఇవేవీ నమ్మకం కలిగించకపోతే, బనాక్--టార్స్కీ వైరుధ్యాభాసాన్ని స్వీకరించడానికి చివరి (బాహ్య) కారణం ఇది. ఇది ప్రపంచంలోనే ఉత్తమమైన గణిత చమత్కారాన్ని వెంటనే ఇస్తుంది:
\begin{enumerate}
	\item[] \emph{ప్రశ్న}. ``Banach--Tarski'' అక్షరాలను పునర్వ్యవస్థీకరిస్తే ఏమవుతుంది?
	\item[] \emph{సమాధానం}. ``Banach--Tarski Banach--Tarski''.
\end{enumerate}

\end{document}
